\documentclass[12pt]{article}
\usepackage[utf8]{inputenc}

\usepackage{../mystyle}

\title{Simons Dissertation Fellowship Proposal}
\author{Joseph Sullivan}
\date{February 2026}

\DeclareMathOperator{\Kom}{Kom}
\DeclareMathOperator{\Cyl}{Cyl}
\DeclareMathOperator{\Stab}{Stab}
\DeclareMathOperator{\Slice}{Slice}
\DeclareMathOperator{\chern}{ch}
\DeclareMathOperator{\todd}{td}

\DeclareMathOperator{\Kos}{Kos}

\DeclareMathOperator{\ext}{ext}
%\DeclareMathOperator{\hom}{hom}

\newcommand{\cQ}{\mathcal{Q}}
\newcommand{\cZ}{\mathcal{Z}}
\newcommand{\LCom}{\mathcal{LC}}

\newcommand{\rddots}{\text{\reflectbox{$\ddots$}}}


\begin{document}

\maketitle


\textbf{Introduction}

My research revolves around \textit{derived categories} in algebraic geometry. Broadly speaking, the derived category of a space $X$ can be considered a more flexible version of the category of vector bundles/coherent sheaves on $X$, but it should also be thought of as a ``higher invariant''\cite{Keller_2007}, capturing data such as Hochshild cohomology, K-theory, and some Hodge theory. My reasons for interest in the derived category are the following:
\begin{itemize}
    \item The existence of \textit{semiorthogonal decompositions}, which describe how one derived category is assembled from smaller pieces.
    \item The existence of interesting derived equivalences between varieties, as in the study of \textit{Fourier-Mukai transforms}, and derived equivalences between varieties and noncommutative algebras (like quiver algebras), as in the study of \textit{tilting bundles} and generalizations.
    \item The objects in a derived category having good moduli theory, with the notion of moduli spaces of \textit{Bridgeland stable} objects, where the spaces vary as a \textit{stability condition} crosses walls in a certain complex manifold.
\end{itemize}

In particular, I like the following program: first understand a derived category by producing semiorthogonal decompositions or equivalences to ``easier'' triangulated categories; and next, use this understanding to produce t-structures on your derived category, which one hopes can be refined to Bridgeland stability conditions.

As far as the immediate future goes, I am currently working on a project with my advisor Aaron Bertram, and my peers Jonathon Fleck and Liebo Pan, about applying Bridgeland stability to find generalizations (for surfaces) of certain theorems controlling the syzygies of a curve embedded by an adjoint divisor.

\textbf{Preparation}

In preparation for this program, I have spent the much of my time following the semiorthogonal decompositions of well-known spaces. Starting with Beilinson's seminal work \cite{Beilinson_1978} producing a full strong exceptional collection on $\P^n$ (the best type of semiorthogonal decomposition), I then moved onto generalizations of Beilinson's argument as in Kapranov's \cite{Kapranov_1988} use of Koszul duality \cite{Beilinson_Ginzburg_Soergel_1996} and Borel-Weil-Bott from representation theory to construct exceptional collections on quadrics. In turn, Kuznetsov \cite{Kuznetsov_2005} repeated Kapranov's argument in families to produce semiorthogonal decompositions for quadric fibrations. Continuing in this direction, I want to better understand semiorthogonal decompositions created by birational geometry as pioneered by Orlov \cite{Orlov_1993} and applied to great success for example in \cite{Tevelev_Torres_2023,Tevelev_2023} producing semiorthogonal decompositions for moduli spaces of vector bundles on curves.

With the structure of these derived categories better understood, I have also tracked how to produce t-structures and their refinements, Bridgeland stability conditions. In particular, projective space and quadrics, both having full strong exceptional collections, admit derived equivalences to quivers via tilting theory \cite{Craw,King}. Transferring the standard heart of a quiver through the equivalence, one can then adapt the theory of GIT quiver moduli \cite{King_1994} as in \cite{Arcara_Bertram_Coskun_Huizenga_2012} to produce so-called \textit{algebraic} Bridgeland stability conditions. This technique has been helpful for producing Bridgeland stability conditions, especially on higher dimensional varieties. On the other hand, Bertram and others have produced so-called geometric stability conditions on \textit{all} surfaces by using torsion pairs and the Bogomolov inequality \cite{Arcara_Bertram_Lieblich_2007}. The project I am working on with Bertram and my peers Fleck and Pan uses these stability conditions.


\textbf{Research Plan}

{\color{red} todo}


\textbf{Career Goals}

I am very much enjoying academia as a PhD candidate, with respect to research, organizing seminars, and being an instructor. I plan to continue onwards to become a postdoc and eventually a professor.


\bibliographystyle{plain} % We choose the "plain" reference style
\bibliography{proposal_bibliography}

\end{document}